\documentclass[11pt,a4paper]{article}
\usepackage[utf8]{inputenc}
\usepackage[T1]{fontenc}
\usepackage[spanish,es-minimal,es-noshorthands]{babel}
\usepackage{lmodern}
\usepackage[sfdefault,lining]{FiraSans}
\usepackage[a4paper,top=2.3cm,bottom=2.2cm,left=1.8cm,right=1.8cm,headheight=1.3cm,headsep=0.4cm,footskip=1cm]{geometry}
\usepackage{amsmath,amssymb}
\usepackage{icomma}
\usepackage{graphicx}
\usepackage[table]{xcolor}
\usepackage{tikz}
\usepackage{fontawesome5}
\usepackage{ccicons}
\usepackage{fancyhdr}
\usepackage{lastpage}
\usepackage{tasks}
\usepackage{needspace}
\usepackage[most]{tcolorbox}
\usepackage{array}
\usepackage[hidelinks]{hyperref}
\AtBeginDocument{\hypersetup{pdftitle={Ficha de ejercicios de la unidad 4: Sistemas de ecuaciones e inecuaciones. 4º ESO},pdfsubject={Matemáticas B 4º ESO},pdfkeywords={sistemas de ecuaciones, método gráfico, sustitución, igualación, reducción, sistemas no lineales, problemas, sistemas de inecuaciones, semiplanos}}}

\definecolor{navy}{HTML}{1B2748}
\definecolor{accent}{HTML}{E8603A}
\definecolor{teal}{HTML}{1F9E90}
\definecolor{soft}{HTML}{EEF1F8}
\definecolor{muted}{HTML}{5D6685}

\setlength{\parindent}{0pt}
\setlength{\parskip}{3pt}
\settasks{label=\alph*), label-format=\bfseries\color{navy}, label-width=1.3em, item-indent=1.8em, column-sep=0.6em, after-item-skip=5pt, before-skip=3pt, after-skip=4pt}

% ---------- capçalera i peu ----------
\newcommand{\llicencia}{\href{https://creativecommons.org/licenses/by-nc-sa/4.0/deed.es}{\ccbyncsa}\ \ Licencia Creative Commons Atribución-NoComercial-CompartirIgual 4.0 Internacional (CC BY-NC-SA 4.0)}
\pagestyle{fancy}
\fancyhf{}
\fancyhead[L]{\raisebox{-0.35\height}{\includegraphics[height=0.95cm]{logo-gva.png}}}
\fancyhead[C]{\small\color{muted}Ficha de ejercicios · Unidad 4: Sistemas de ecuaciones e inecuaciones · 4\textsuperscript{o} ESO B}
\fancyhead[R]{\raisebox{-0.35\height}{\includegraphics[height=0.95cm]{logo-vermellar.png}}}
\fancyfoot[L]{\scriptsize\color{muted}\llicencia}
\fancyfoot[R]{\small\color{muted}\thepage}
\renewcommand{\headrulewidth}{0.4pt}
\renewcommand{\headrule}{\color{muted!40}\hrule height \headrulewidth}
\fancypagestyle{primera}{\fancyhf{}\renewcommand{\headrulewidth}{0pt}
  \fancyfoot[L]{\scriptsize\color{muted}\llicencia}
  \fancyfoot[R]{\small\color{muted}\thepage}}

% ---------- nivells i exercicis ----------
\newcommand{\estrelles}[1]{{\color{accent}\ifcase#1\or\faStar\faStar[regular]\faStar[regular]\or\faStar\faStar\faStar[regular]\or\faStar\faStar\faStar\fi}}
\newcounter{exer}
\newcommand{\ex}[1]{\par\needspace{4\baselineskip}\medskip\refstepcounter{exer}\noindent{\color{navy}\bfseries\theexer.}\ \ {\small\estrelles{#1}}\ \ }

\newcommand{\seccio}[2]{\par\needspace{16\baselineskip}\bigskip
  \begin{tcolorbox}[colback=navy,colframe=navy,coltext=white,arc=2mm,boxrule=0pt,left=3mm,right=3mm,top=1.2mm,bottom=1.2mm]
  {\large\bfseries #1\quad #2}
  \end{tcolorbox}\vspace{-2pt}}

\newcommand{\exf}[1]{\par\medskip\refstepcounter{exer}\noindent{\color{navy}\bfseries\theexer.}\ \ {\small\estrelles{#1}}\ \ }
\newcolumntype{C}[1]{>{\centering\arraybackslash}p{#1}}
% quadre de recordatori (només a la fitxa d'exercicis)
\newcommand{\recorda}[1]{\par\smallskip\begin{tcolorbox}[colback=accent!7,colframe=accent!55,arc=2mm,boxrule=0.5pt,left=2.5mm,right=2.5mm,top=1mm,bottom=1mm]\small{\color{accent}\faLightbulb[regular]}\ \textbf{Recuerda:}\ #1\end{tcolorbox}}
\usetikzlibrary{calc}
\newcommand{\res}[1]{#1}\newcommand{\rest}[1]{#1}

\def\fmneed{11}% que el títol d'una secció no quede sol al final del full
% =====================================================================
% Capçalera FentMates de les fitxes: un color per a cada curs.
% Els colors són els de la pàgina inicial (LEVELS a assets/js/materials.js);
% el logo ve de _eines/logo/pdf/ (python3 _eines/logo/genera.py).
%
% Ús, al final del preàmbul de la unitat:
%   % =====================================================================
% Capçalera FentMates de les fitxes: un color per a cada curs.
% Els colors són els de la pàgina inicial (LEVELS a assets/js/materials.js);
% el logo ve de _eines/logo/pdf/ (python3 _eines/logo/genera.py).
%
% Ús, al final del preàmbul de la unitat:
%   % =====================================================================
% Capçalera FentMates de les fitxes: un color per a cada curs.
% Els colors són els de la pàgina inicial (LEVELS a assets/js/materials.js);
% el logo ve de _eines/logo/pdf/ (python3 _eines/logo/genera.py).
%
% Ús, al final del preàmbul de la unitat:
%   \input{../comu/fentmates.tex}
%   \fmcurs{1eso}
% i, al començament del document:
%   \capcalera{Matemàtiques}{1r ESO}{Fitxa d'exercicis · Unitat 2}{Divisibilitat}
% Opcional abans de \input: \def\fmneed{14} (espai mínim abans de cada secció).
% =====================================================================
\makeatletter
\def\fm@logodir{../../logo/pdf}
\providecommand{\fmneed}{10}
\def\fm@cursnom{}

% clau del curs, color i color del text damunt del color
\newcommand{\fm@curs}[3]{\@namedef{fm@color@#1}{#2}\@namedef{fm@text@#1}{#3}}
\fm@curs{1eso}{E6B84A}{navy}
\fm@curs{2eso}{3B82F6}{white}
\fm@curs{3eso}{FF7A50}{white}
\fm@curs{4eso-a}{A78BFA}{white}
\fm@curs{4eso-b}{6D28D9}{white}
\fm@curs{1batx-mi}{9B1C1C}{white}
\fm@curs{1batx-ccss}{EF6B6B}{white}
\fm@curs{2batx-mii}{1E6B47}{white}
\fm@curs{2batx-ccss}{2BB3A3}{white}

\newcommand{\fmcurs}[1]{%
  \@ifundefined{fm@color@#1}{\PackageError{fentmates}{Curs desconegut: #1}{}}{}%
  \edef\fm@tmp{\noexpand\definecolor{curs}{HTML}{\@nameuse{fm@color@#1}}%
               \noexpand\colorlet{curstext}{\@nameuse{fm@text@#1}}}\fm@tmp
  \colorlet{cursfosc}{curs!72!black}%
  \def\fm@logo{\fm@logodir/vertical-#1.pdf}%
  \def\fm@marca{\fm@logodir/marca-#1.pdf}}

% ---------- capçalera de la primera pàgina ----------
% #1 matèria (petit), #2 curs, #3 tipus de document i unitat, #4 títol
\newcommand{\capcalera}[4]{%
  \gdef\fm@cursnom{#2}%
  \par\noindent
  \begin{tcolorbox}[colback=curs,colframe=curs,boxrule=0pt,arc=2.2mm,width=3.5cm,height=2.15cm,
      valign=center,left=3mm,right=2mm,top=1mm,bottom=1mm,nobeforeafter,box align=center]
    \color{curstext}{\footnotesize #1}\par\vspace{1pt}{\fontsize{19}{21}\selectfont\bfseries #2\par}%
  \end{tcolorbox}\hspace{5mm}%
  \begin{minipage}[c]{\dimexpr\linewidth-3.5cm-5mm-2.9cm\relax}
    {\color{cursfosc}\small\bfseries\MakeUppercase{#3}}\par\vspace{2pt}
    {\color{navy}\fontsize{20}{23}\selectfont\bfseries #4\par}
  \end{minipage}\hfill
  \begin{minipage}[c]{2.6cm}\raggedleft\includegraphics[width=2.6cm]{\fm@logo}\end{minipage}\par
  \vspace{2pt}}

% ---------- barres de secció del color del curs ----------
\renewcommand{\seccio}[2]{\par\needspace{\fmneed\baselineskip}\bigskip
  \begin{tcolorbox}[colback=curs,colframe=curs,coltext=curstext,arc=2mm,boxrule=0pt,left=3mm,right=3mm,top=1.2mm,bottom=1.2mm]
  {\large\bfseries #1\quad #2}
  \end{tcolorbox}\vspace{-2pt}}

% ---------- estrelles de dificultat del color del curs (a Batxillerat no se'n posen) ----------
\@ifundefined{estrelles}{}{%
  \renewcommand{\estrelles}[1]{{\color{curs}\ifcase#1\or\faStar\faStar[regular]\faStar[regular]\or\faStar\faStar\faStar[regular]\or\faStar\faStar\faStar\fi}}}

% ---------- capçalera i peu de les pàgines ----------
% primera pàgina: logos institucionals al peu
\fancypagestyle{primera}{\fancyhf{}\renewcommand{\headrulewidth}{0pt}
  \fancyfoot[L]{\raisebox{-0.32\height}{\includegraphics[height=0.78cm]{logo-gva.png}}\hspace{5mm}%
    \raisebox{-0.32\height}{\includegraphics[height=0.78cm]{logo-vermellar.png}}}
  \fancyfoot[R]{\scriptsize\color{muted}\href{https://creativecommons.org/licenses/by-nc-sa/4.0/deed.ca}{\ccbyncsa}\ \ CC BY-NC-SA 4.0 · fentmates.org\qquad\small\thepage}}
% la resta: el símbol a l'esquerra i el curs a la dreta, amb el títol al centre com abans
\fancyhead[L]{\raisebox{-0.3\height}{\includegraphics[height=0.62cm]{\fm@marca}}}
\fancyhead[R]{\small\bfseries\color{cursfosc}\fm@cursnom}
\renewcommand{\headrule}{\color{curs!45}\hrule height \headrulewidth}
\makeatother

%   \fmcurs{1eso}
% i, al començament del document:
%   \capcalera{Matemàtiques}{1r ESO}{Fitxa d'exercicis · Unitat 2}{Divisibilitat}
% Opcional abans de \input: \def\fmneed{14} (espai mínim abans de cada secció).
% =====================================================================
\makeatletter
\def\fm@logodir{../../logo/pdf}
\providecommand{\fmneed}{10}
\def\fm@cursnom{}

% clau del curs, color i color del text damunt del color
\newcommand{\fm@curs}[3]{\@namedef{fm@color@#1}{#2}\@namedef{fm@text@#1}{#3}}
\fm@curs{1eso}{E6B84A}{navy}
\fm@curs{2eso}{3B82F6}{white}
\fm@curs{3eso}{FF7A50}{white}
\fm@curs{4eso-a}{A78BFA}{white}
\fm@curs{4eso-b}{6D28D9}{white}
\fm@curs{1batx-mi}{9B1C1C}{white}
\fm@curs{1batx-ccss}{EF6B6B}{white}
\fm@curs{2batx-mii}{1E6B47}{white}
\fm@curs{2batx-ccss}{2BB3A3}{white}

\newcommand{\fmcurs}[1]{%
  \@ifundefined{fm@color@#1}{\PackageError{fentmates}{Curs desconegut: #1}{}}{}%
  \edef\fm@tmp{\noexpand\definecolor{curs}{HTML}{\@nameuse{fm@color@#1}}%
               \noexpand\colorlet{curstext}{\@nameuse{fm@text@#1}}}\fm@tmp
  \colorlet{cursfosc}{curs!72!black}%
  \def\fm@logo{\fm@logodir/vertical-#1.pdf}%
  \def\fm@marca{\fm@logodir/marca-#1.pdf}}

% ---------- capçalera de la primera pàgina ----------
% #1 matèria (petit), #2 curs, #3 tipus de document i unitat, #4 títol
\newcommand{\capcalera}[4]{%
  \gdef\fm@cursnom{#2}%
  \par\noindent
  \begin{tcolorbox}[colback=curs,colframe=curs,boxrule=0pt,arc=2.2mm,width=3.5cm,height=2.15cm,
      valign=center,left=3mm,right=2mm,top=1mm,bottom=1mm,nobeforeafter,box align=center]
    \color{curstext}{\footnotesize #1}\par\vspace{1pt}{\fontsize{19}{21}\selectfont\bfseries #2\par}%
  \end{tcolorbox}\hspace{5mm}%
  \begin{minipage}[c]{\dimexpr\linewidth-3.5cm-5mm-2.9cm\relax}
    {\color{cursfosc}\small\bfseries\MakeUppercase{#3}}\par\vspace{2pt}
    {\color{navy}\fontsize{20}{23}\selectfont\bfseries #4\par}
  \end{minipage}\hfill
  \begin{minipage}[c]{2.6cm}\raggedleft\includegraphics[width=2.6cm]{\fm@logo}\end{minipage}\par
  \vspace{2pt}}

% ---------- barres de secció del color del curs ----------
\renewcommand{\seccio}[2]{\par\needspace{\fmneed\baselineskip}\bigskip
  \begin{tcolorbox}[colback=curs,colframe=curs,coltext=curstext,arc=2mm,boxrule=0pt,left=3mm,right=3mm,top=1.2mm,bottom=1.2mm]
  {\large\bfseries #1\quad #2}
  \end{tcolorbox}\vspace{-2pt}}

% ---------- estrelles de dificultat del color del curs (a Batxillerat no se'n posen) ----------
\@ifundefined{estrelles}{}{%
  \renewcommand{\estrelles}[1]{{\color{curs}\ifcase#1\or\faStar\faStar[regular]\faStar[regular]\or\faStar\faStar\faStar[regular]\or\faStar\faStar\faStar\fi}}}

% ---------- capçalera i peu de les pàgines ----------
% primera pàgina: logos institucionals al peu
\fancypagestyle{primera}{\fancyhf{}\renewcommand{\headrulewidth}{0pt}
  \fancyfoot[L]{\raisebox{-0.32\height}{\includegraphics[height=0.78cm]{logo-gva.png}}\hspace{5mm}%
    \raisebox{-0.32\height}{\includegraphics[height=0.78cm]{logo-vermellar.png}}}
  \fancyfoot[R]{\scriptsize\color{muted}\href{https://creativecommons.org/licenses/by-nc-sa/4.0/deed.ca}{\ccbyncsa}\ \ CC BY-NC-SA 4.0 · fentmates.org\qquad\small\thepage}}
% la resta: el símbol a l'esquerra i el curs a la dreta, amb el títol al centre com abans
\fancyhead[L]{\raisebox{-0.3\height}{\includegraphics[height=0.62cm]{\fm@marca}}}
\fancyhead[R]{\small\bfseries\color{cursfosc}\fm@cursnom}
\renewcommand{\headrule}{\color{curs!45}\hrule height \headrulewidth}
\makeatother

%   \fmcurs{1eso}
% i, al començament del document:
%   \capcalera{Matemàtiques}{1r ESO}{Fitxa d'exercicis · Unitat 2}{Divisibilitat}
% Opcional abans de \input: \def\fmneed{14} (espai mínim abans de cada secció).
% =====================================================================
\makeatletter
\def\fm@logodir{../../logo/pdf}
\providecommand{\fmneed}{10}
\def\fm@cursnom{}

% clau del curs, color i color del text damunt del color
\newcommand{\fm@curs}[3]{\@namedef{fm@color@#1}{#2}\@namedef{fm@text@#1}{#3}}
\fm@curs{1eso}{E6B84A}{navy}
\fm@curs{2eso}{3B82F6}{white}
\fm@curs{3eso}{FF7A50}{white}
\fm@curs{4eso-a}{A78BFA}{white}
\fm@curs{4eso-b}{6D28D9}{white}
\fm@curs{1batx-mi}{9B1C1C}{white}
\fm@curs{1batx-ccss}{EF6B6B}{white}
\fm@curs{2batx-mii}{1E6B47}{white}
\fm@curs{2batx-ccss}{2BB3A3}{white}

\newcommand{\fmcurs}[1]{%
  \@ifundefined{fm@color@#1}{\PackageError{fentmates}{Curs desconegut: #1}{}}{}%
  \edef\fm@tmp{\noexpand\definecolor{curs}{HTML}{\@nameuse{fm@color@#1}}%
               \noexpand\colorlet{curstext}{\@nameuse{fm@text@#1}}}\fm@tmp
  \colorlet{cursfosc}{curs!72!black}%
  \def\fm@logo{\fm@logodir/vertical-#1.pdf}%
  \def\fm@marca{\fm@logodir/marca-#1.pdf}}

% ---------- capçalera de la primera pàgina ----------
% #1 matèria (petit), #2 curs, #3 tipus de document i unitat, #4 títol
\newcommand{\capcalera}[4]{%
  \gdef\fm@cursnom{#2}%
  \par\noindent
  \begin{tcolorbox}[colback=curs,colframe=curs,boxrule=0pt,arc=2.2mm,width=3.5cm,height=2.15cm,
      valign=center,left=3mm,right=2mm,top=1mm,bottom=1mm,nobeforeafter,box align=center]
    \color{curstext}{\footnotesize #1}\par\vspace{1pt}{\fontsize{19}{21}\selectfont\bfseries #2\par}%
  \end{tcolorbox}\hspace{5mm}%
  \begin{minipage}[c]{\dimexpr\linewidth-3.5cm-5mm-2.9cm\relax}
    {\color{cursfosc}\small\bfseries\MakeUppercase{#3}}\par\vspace{2pt}
    {\color{navy}\fontsize{20}{23}\selectfont\bfseries #4\par}
  \end{minipage}\hfill
  \begin{minipage}[c]{2.6cm}\raggedleft\includegraphics[width=2.6cm]{\fm@logo}\end{minipage}\par
  \vspace{2pt}}

% ---------- barres de secció del color del curs ----------
\renewcommand{\seccio}[2]{\par\needspace{\fmneed\baselineskip}\bigskip
  \begin{tcolorbox}[colback=curs,colframe=curs,coltext=curstext,arc=2mm,boxrule=0pt,left=3mm,right=3mm,top=1.2mm,bottom=1.2mm]
  {\large\bfseries #1\quad #2}
  \end{tcolorbox}\vspace{-2pt}}

% ---------- estrelles de dificultat del color del curs (a Batxillerat no se'n posen) ----------
\@ifundefined{estrelles}{}{%
  \renewcommand{\estrelles}[1]{{\color{curs}\ifcase#1\or\faStar\faStar[regular]\faStar[regular]\or\faStar\faStar\faStar[regular]\or\faStar\faStar\faStar\fi}}}

% ---------- capçalera i peu de les pàgines ----------
% primera pàgina: logos institucionals al peu
\fancypagestyle{primera}{\fancyhf{}\renewcommand{\headrulewidth}{0pt}
  \fancyfoot[L]{\raisebox{-0.32\height}{\includegraphics[height=0.78cm]{logo-gva.png}}\hspace{5mm}%
    \raisebox{-0.32\height}{\includegraphics[height=0.78cm]{logo-vermellar.png}}}
  \fancyfoot[R]{\scriptsize\color{muted}\href{https://creativecommons.org/licenses/by-nc-sa/4.0/deed.ca}{\ccbyncsa}\ \ CC BY-NC-SA 4.0 · fentmates.org\qquad\small\thepage}}
% la resta: el símbol a l'esquerra i el curs a la dreta, amb el títol al centre com abans
\fancyhead[L]{\raisebox{-0.3\height}{\includegraphics[height=0.62cm]{\fm@marca}}}
\fancyhead[R]{\small\bfseries\color{cursfosc}\fm@cursnom}
\renewcommand{\headrule}{\color{curs!45}\hrule height \headrulewidth}
\makeatother

\fmcurs{4eso-b}

\begin{document}
\thispagestyle{primera}

\capcalera{Matemáticas B}{4\textsuperscript{o} ESO}{Ficha de ejercicios · Unidad 4}{\raggedright Sistemas de ecuaciones e inecuaciones}
\vspace{10pt}
Nombre y apellidos: \rule{6.6cm}{0.4pt}\hfill Grupo: \rule{1.6cm}{0.4pt}\hfill Fecha: \rule{2.2cm}{0.4pt}
\vspace{8pt}

\begin{tcolorbox}[colback=soft,colframe=soft,arc=2mm,left=3mm,right=3mm,top=1.5mm,bottom=1.5mm]
\small Ejercicios graduados por niveles de dificultad:\quad \estrelles{1}\ básico \qquad \estrelles{2}\ medio \qquad \estrelles{3}\ avanzado.\\[2pt] Resuélvelos en el cuaderno mostrando todos los pasos. Comprueba las soluciones de los sistemas en las dos ecuaciones y, en los problemas, responde con una frase.
\end{tcolorbox}

% =====================================================
\seccio{1}{Sistemas lineales. Método gráfico}

\exf{1} Comprueba si $(2,-1)$ y $(1,1)$ son solución del sistema $\begin{cases}3x+2y=4\\[2pt] x-y=3\end{cases}$.

\ex{2} Resuelve gráficamente los sistemas:\begin{tasks}[after-item-skip=9pt](3)
  \task $\begin{cases}x+y=5\\[2pt] x-y=1\end{cases}$
  \task $\begin{cases}2x-y=1\\[2pt] x+y=5\end{cases}$
  \task $\begin{cases}x-2y=-4\\[2pt] x+y=2\end{cases}$
\end{tasks}\recorda{Haz una tabla de valores con dos puntos de cada recta. La solución es el punto de corte.}

\ex{2} Clasifica los sistemas sin resolverlos:\begin{tasks}[after-item-skip=9pt](3)
  \task $\begin{cases}2x+3y=5\\[2pt] 4x+6y=10\end{cases}$
  \task $\begin{cases}x-y=2\\[2pt] 3x-3y=1\end{cases}$
  \task $\begin{cases}x+2y=3\\[2pt] 2x-y=1\end{cases}$
\end{tasks}

\ex{3} Halla el valor de $k$ para que el sistema $\begin{cases}kx+2y=4\\[2pt] 3x+y=5\end{cases}$ sea incompatible.


% =====================================================
\seccio{2}{Métodos de sustitución, igualación y reducción}

\exf{2} Resuelve por el método de sustitución:\begin{tasks}[after-item-skip=16pt](2)
  \task $\begin{cases}y=2x-3\\[2pt] 3x+y=7\end{cases}$
  \task $\begin{cases}x+3y=10\\[2pt] 2x-y=6\end{cases}$
  \task $\begin{cases}\dfrac x2+y=5\\[9pt] 3(x-1)-2y=3\end{cases}$
  \task $\begin{cases}\dfrac{x+1}{5}=y-1\\[9pt] 2(x-y)+y=6\end{cases}$
  \task $\begin{cases}\dfrac x3-y=4\\[9pt] 2(x+y)-3x=-10\end{cases}$
\end{tasks}

\ex{2} Resuelve por el método de igualación:\begin{tasks}[after-item-skip=16pt](2)
  \task $\begin{cases}y=x+4\\[2pt] y=3x\end{cases}$
  \task $\begin{cases}x=2y-1\\[2pt] x=8-y\end{cases}$
  \task $\begin{cases}x=\dfrac{y+7}{3}\\[9pt] x=\dfrac{11-y}{3}\end{cases}$
  \task $\begin{cases}\dfrac y2=\dfrac{x-1}{3}\\[9pt] y=6-x\end{cases}$
  \task $\begin{cases}x=\dfrac{y-6}{2}\\[9pt] x+y=3\end{cases}$
\end{tasks}

\ex{2} Resuelve por el método de reducción:\begin{tasks}[after-item-skip=16pt](2)
  \task $\begin{cases}2x+3y=12\\[2pt] 2x-y=4\end{cases}$
  \task $\begin{cases}3x+2y=7\\[2pt] 5x-3y=18\end{cases}$
  \task $\begin{cases}4x+3y=1\\[2pt] 3x-2y=5\end{cases}$
  \task $\begin{cases}\dfrac{x+y}{2}+\dfrac{x-y}{4}=\dfrac52\\[9pt] 2(x+1)-3(y-1)=8\end{cases}$
  \task $\begin{cases}\dfrac{x-1}{3}+\dfrac{y+3}{2}=\dfrac43\\[9pt] 3(x+y)-2(x-y)=-3\end{cases}$
  \task $\begin{cases}\dfrac x5-\dfrac y2=0\\[9pt] \dfrac{x+1}{3}+\dfrac{y-2}{4}=2\end{cases}$
\end{tasks}\recorda{Multiplica las ecuaciones para que una incógnita tenga el mismo coeficiente (o el opuesto) y réstalas (o súmalas).}

\ex{2} Resuelve el sistema $\begin{cases}x+y=9\\[2pt] x-y=3\end{cases}$ por los tres métodos algebraicos y comprueba que obtienes la misma solución.

\ex{2} Simplifica las ecuaciones y resuelve:\begin{tasks}[after-item-skip=16pt](3)
  \task $\begin{cases}\dfrac{x+1}{2}+\dfrac y3=2\\[9pt] 2(x-y)=x-3\end{cases}$
  \task $\begin{cases}\dfrac x3+\dfrac y2=4\\[9pt] \dfrac x2-\dfrac y4=2\end{cases}$
  \task $\begin{cases}3(x-1)+y=6\\[2pt] x-2(y+1)=-6\end{cases}$
\end{tasks}

\ex{2} Resuelve y clasifica:\begin{tasks}[after-item-skip=9pt](2)
  \task $\begin{cases}2x-4y=6\\[2pt] -x+2y=-3\end{cases}$
  \task $\begin{cases}x+y=3\\[2pt] 2x+2y=7\end{cases}$
\end{tasks}


% =====================================================
\seccio{3}{Sistemas no lineales}

\exf{2} Resuelve por el método de sustitución:\begin{tasks}[after-item-skip=9pt](3)
  \task $\begin{cases}x+y=5\\[2pt] x^2+y^2=13\end{cases}$
  \task $\begin{cases}x-y=2\\[2pt] xy=15\end{cases}$
  \task $\begin{cases}y=x^2-4\\[2pt] y=2x-1\end{cases}$
\end{tasks}

\ex{2} Resuelve por el método de reducción:\begin{tasks}[after-item-skip=9pt](2)
  \task $\begin{cases}x^2+y^2=25\\[2pt] x^2-y^2=7\end{cases}$
  \task $\begin{cases}2x^2+y^2=6\\[2pt] x^2-y^2=3\end{cases}$
\end{tasks}

\ex{3} Resuelve:\begin{tasks}[after-item-skip=16pt](2)
  \task $\begin{cases}x+y=1\\[9pt] \dfrac1x+\dfrac1y=-\dfrac16\end{cases}$
  \task $\begin{cases}x^2+y^2=10\\[2pt] xy=3\end{cases}$
\end{tasks}

\ex{2} Resuelve:\begin{tasks}[after-item-skip=16pt](2)
  \task $\begin{cases}x^2+y^2=10\\[2pt] x+y=4\end{cases}$
  \task $\begin{cases}y=x^2-2x\\[2pt] y=x+4\end{cases}$
  \task $\begin{cases}\dfrac1x+\dfrac1y=\dfrac56\\[9pt] x+y=5\end{cases}$
  \task $\begin{cases}xy=12\\[2pt] x^2+y^2=25\end{cases}$
\end{tasks}

\ex{3} Resuelve los sistemas siguientes, en los que hay raíces. Comprueba las soluciones:\begin{tasks}[after-item-skip=9pt](3)
  \task $\begin{cases}\sqrt x+y=5\\[2pt] x-y=1\end{cases}$
  \task $\begin{cases}\sqrt{x+y}=3\\[2pt] x-y=1\end{cases}$
  \task $\begin{cases}\sqrt{x+3}=y\\[2pt] x+y=9\end{cases}$
\end{tasks}\recorda{Aísla la raíz y eleva al cuadrado; después comprueba cada solución en el sistema inicial, porque al elevar al cuadrado pueden aparecer soluciones falsas.}

\ex{2} Sin dibujar, averigua si la recta $y=x+3$ corta a la circunferencia $x^2+y^2=4$.


% =====================================================
\seccio{4}{Problemas con sistemas}

\exf{1} En una clase hay 28 alumnos, y el número de chicas supera en 4 al de chicos. ¿Cuántas chicas y cuántos chicos hay?

\ex{2} Tres entradas de adulto y dos de niño cuestan 33 €, y dos de adulto y tres de niño, 27 €. ¿Cuánto cuesta cada entrada?

\ex{2} Las cifras de un número de dos cifras suman 11. Si las intercambiamos, el número disminuye en 27. ¿Cuál es el número?

\ex{2} Queremos hacer 20 kg de una mezcla de café a 9,40 €/kg con café de 8 €/kg y café de 12 €/kg. ¿Cuántos kilos debemos poner de cada clase?

\ex{2} El perímetro de un rectángulo es de 28 m y la diagonal mide 10 m. ¿Cuánto miden los lados?

\ex{3} La hipotenusa de un triángulo rectángulo mide 13 cm y el área, 30 cm$^2$. ¿Cuánto miden los catetos?


% =====================================================
\seccio{5}{Sistemas de inecuaciones con una incógnita}

\exf{1} Resuelve los sistemas y representa la solución:\begin{tasks}[after-item-skip=9pt](2)
  \task $\begin{cases}x-3<0\\[2pt] x+1\ge0\end{cases}$
  \task $\begin{cases}2x+1>5\\[2pt] 3x-2\le10\end{cases}$
  \task $\begin{cases}x+4<1\\[2pt] 2x>-2\end{cases}$
  \task $\begin{cases}5-x\le2\\[2pt] 4x+1\ge9\end{cases}$
\end{tasks}\recorda{Resuelve cada inecuación y haz la intersección de las soluciones: dibújalas una encima de la otra.}

\ex{2} Resuelve:\begin{tasks}[after-item-skip=16pt](2)
  \task $\begin{cases}3(x-1)\ge x+1\\[9pt] \dfrac{x+2}{3}<2\end{cases}$
  \task $\begin{cases}2x-5<x+1\\[9pt] \dfrac x2+1>0\end{cases}$
\end{tasks}

\ex{3} Resuelve (alguna inecuación es de segundo grado):\begin{tasks}[after-item-skip=9pt](2)
  \task $\begin{cases}x^2-9\le0\\[2pt] 2x-1>0\end{cases}$
  \task $\begin{cases}x^2-4x>0\\[2pt] x+2\ge0\end{cases}$
\end{tasks}

\ex{2} Un taxi cobra 3 € por la bajada de bandera y 1,20 € por kilómetro. Si queremos gastar como mínimo 9 € y como máximo 15 €, ¿cuántos kilómetros podemos hacer?


% =====================================================
\seccio{6}{Inecuaciones con dos incógnitas}

\exf{1} Representa la solución de las inecuaciones:\begin{tasks}(4)
  \task $x+y\le3$
  \task $2x-y>4$
  \task $y\ge2x$
  \task $x<-1$
\end{tasks}\recorda{Dibuja la recta frontera (continua si hay $\le$ o $\ge$, discontinua si hay $<$ o $>$) y prueba un punto que no esté en ella.}

\ex{1} Indica cuáles de estos puntos son solución de $2x+3y\le6$: $A(1,1)$, $B(3,-2)$, $C(0,5)$.


% =====================================================
\seccio{7}{Sistemas de inecuaciones con dos incógnitas}

\exf{2} Representa la región solución de cada sistema y halla sus vértices:\begin{tasks}[after-item-skip=9pt](3)
  \task $\begin{cases}x+y\le5\\[2pt] x\ge0\\[2pt] y\ge0\end{cases}$
  \task $\begin{cases}y\ge x\\[2pt] y\le4\\[2pt] x\ge1\end{cases}$
  \task $\begin{cases}x+2y\le8\\[2pt] x-y\ge-1\\[2pt] y\ge0\end{cases}$
\end{tasks}\recorda{La solución es la intersección de los semiplanos. Los vértices son los puntos de corte de las rectas frontera.}

\ex{3} Representa la región solución de $\begin{cases}x+y\ge2\\[2pt] x\ge0\\[2pt] y\ge0\end{cases}$. ¿Está acotada?

\ex{2} Un taller hace mesas ($x$) y sillas ($y$). Cada día puede hacer como máximo 10 piezas en total, al menos 2 mesas, y nunca más sillas que mesas. Escribe el sistema de inecuaciones, represéntalo y di si pueden hacer 4 mesas y 5 sillas.

\vfill
\begin{center}\color{muted}\small\faCheckCircle[regular]\ \ ¡Ánimo! Comprueba tus resultados con las herramientas interactivas de los apuntes de la unidad.\end{center}

\end{document}
